\section{从代码修复到安全任务：能力迁移与场景重构}

\subsection{安全任务与普通修复任务的结构差异}
安全任务不止 ``让测试通过''，还要求回答三个更难的问题：漏洞是否真实可触发、触发条件是什么、是否可跨路径复现。相比一般代码修复，漏洞挖掘更依赖动态行为观测、跨函数数据流推理和反事实验证。

\begin{knowledgebox}{安全 Agent 的最小任务闭环}
\begin{enumerate}
    \item 假设某路径可能存在越界/溢出/未检查条件；
    \item 生成可执行输入或触发脚本；
    \item 在 sanitizer/debugger 下运行并收集证据；
    \item 失败则修正假设并继续探索，直到复现或证伪。
\end{enumerate}
\end{knowledgebox}

\subsection{CTF/Enigma 作为过渡台阶}
讲座给出 CTF 与 Enigma 类任务作为中间层：它们比 IDE 内日常补丁更强调利用链和触发链，但比真实大型系统漏洞仍可控。这个台阶有价值，因为它让团队能在可评测环境里迭代工具与控制流。

\begin{importantbox}{为什么先做 CTF 再做真实漏洞}
CTF 允许在可重复环境里快速验证 ``工具是否有用''；真实漏洞则更多受上下文噪声、依赖环境、输入空间复杂度影响。先把 Agent 基础循环打磨好，再向真实系统迁移，成功率更高。
\end{importantbox}

\subsection{动态分析工具进入主流程}
在 00:50 后半段，主讲人专门讲了 sanitizer 作为动态分析器的角色。与纯静态判断不同，sanitizer 在运行时直接提供 ``是否触发非法内存行为'' 的信号，这使它天然适合成为 Agent 的 oracle。

\begin{warningbox}{仅依赖静态分类器会陷入“看起来像漏洞”}
很多函数局部看似危险，但全局调用链可能已做边界检查；反之，局部看似安全，跨模块组合后仍可能可利用。缺少动态验证时，误报与漏报都会显著增加。
\end{warningbox}

\subsection{本章小结}
安全任务不是代码修复任务的简单放大，而是目标函数变化。只追求测试通过无法回答 ``是否真有可利用漏洞''，必须引入动态执行证据和更强的验证回路。

